\documentclass[11pt,a4paper]{article}

\usepackage[T1]{fontenc}
\usepackage[utf8]{inputenc}
\usepackage[margin=1.1in]{geometry}
\usepackage{mathpazo}
\usepackage{microtype}
\usepackage{enumitem}
\usepackage{parskip}
\usepackage{titling}
\usepackage[hidelinks]{hyperref}

\setlength{\parindent}{0pt}
\setlength{\parskip}{6pt}

\setlist[itemize]{leftmargin=1.4em, topsep=2pt, itemsep=2pt}
\setlist[enumerate]{leftmargin=1.6em, topsep=2pt, itemsep=2pt}

\setlength{\droptitle}{-3em}

\title{When Your Legal AI Is Wrong: A Diagnostic Ladder for
Lawyers\\[0.3em]
\large Five questions --- structure, meaning, context,
groundedness, environment --- for the moment when an AI-assisted
output looks plausible but is not quite right.}

\author{Lodewijk Bonebakker\thanks{Drafting assistance: Claude
Opus 4.7 (Anthropic). The author chose the framework, scope, and
voice; the model produced prose under direction; the author
edited. The author is not a practicing attorney; nothing in this
article is legal advice.}}

\date{April 2026}

\begin{document}
\maketitle

\noindent
A senior associate at your firm has just sent you a memorandum
drafted with the help of an AI tool. The memo cites the right
statute, paraphrases the right line of cases, and arrives at what
looks like the right conclusion. But something feels off, and when
you read the draft carefully you cannot quite say what. The cases
are real this time. The reasoning is internally consistent. The
output is, in some specific way, wrong, and ``the AI hallucinated''
does not capture what went wrong because the AI, in the narrow
sense, did not.

The fabricated-cases scandal in \emph{Mata v.~Avianca} [1] taught
the profession to fear hallucination. The 2024 Stanford RegLab
study by Magesh and colleagues [2] then taught us that even the
retrieval-grounded legal research tools sold by the major
publishers --- Westlaw AI, Lexis+ AI, and others --- still
hallucinate at meaningful rates, between roughly seventeen and
thirty-three percent of queries in their evaluation. These two
data points have shaped how the profession thinks about AI errors:
hallucination is the headline. But headline failures are also the
easy ones to catch. The harder failures, the ones that survive a
quick cite-check and arrive on a partner's desk looking respectable,
are not hallucinations at all.

\section*{The vocabulary the profession has is too narrow}

The everyday vocabulary lawyers use to describe AI errors collapses
several distinct phenomena into one word. ``It hallucinated.''
``It was wrong.'' ``It missed the point.'' Each of these names a
real failure but does not point at where in the workflow the
failure happened or what would prevent it next time. A lawyer who
discovers a fabricated citation does not respond the same way as a
lawyer who discovers that the memo missed the actual question the
client asked. The profession's existing language treats both as
``the AI got it wrong'' and offers no obvious next step.

What is needed is a small, structured vocabulary --- a diagnostic
ladder a lawyer can walk quickly when a draft arrives and
something feels off. The ladder this article describes has five
rungs: \emph{structure}, \emph{meaning}, \emph{context},
\emph{groundedness}, and \emph{environment}. The rungs run from
what the AI is best at to what only the lawyer can supply. Most
review workflows attend carefully to the lowest rung and to
groundedness; the middle rungs --- meaning and context --- are
checked less consistently, and environment usually not at all.
That asymmetry is where the silent damage tends to live.

The five rungs address single-output failures --- everything
diagnosable when a specific draft is in front of you. Six
additional failure classes exist alongside them that the
rung-by-rung walk cannot catch; they require different detection
and remediation. The article maps both before demonstrating the
ladder in practice.

The examples below draw from the work most law firms now ask AI to
help with: legal research, document review, drafting (memos,
briefs, contracts, client correspondence), analysis (issue
spotting, application of rules to facts), and client
communication. The ladder is the same in each.

\section*{Structure: is the output well-formed?}

A model's output is structurally correct when it satisfies the
formal conventions of the artifact type: a properly Bluebooked
citation, a contract whose section numbering is consistent and
whose defined terms are capitalized in the right places, a
privilege log in the format the court accepts, a filing in the
format the e-filing system requires. Structure is mechanical.
Modern AI tools do it well most of the time, and what they get
wrong is usually catchable by automated checkers --- citation
linters, contract-comparison tools, court-filing validators.

The diagnostic risk on this rung is not that structural failures
are common; it is that they are the easiest to catch, and a
reviewer who confirms structure and stops has confirmed the least
interesting property of the output. A correctly Bluebooked
citation can still cite a case that does not stand for the
proposition the memo says it does. A contract with consistent
defined terms can still indemnify the wrong party. Structure is
necessary, not sufficient.

\section*{Meaning: is the output internally coherent?}

A draft can be structurally fine and internally inconsistent. A
research memo argues both sides of an issue without resolving
them. A contract's recitals describe a transaction whose
operative covenants do not match. An opinion letter's conclusion
does not follow from the analysis preceding it. A drafted
amendment redefines a term in a way the underlying agreement does
not anticipate. A redlined clause subtly inverts the original
allocation of risk. These failures are detectable by careful
reading without recourse to any outside source. They are failures
of the document considered as a self-contained artifact.

The mechanism is the same one that produces fluent legal prose in
the first place. A language model trained to predict the next
token optimizes for the plausibility of each token in isolation.
Plausibility is local. A clause can be plausible in isolation and
still be inconsistent with one three pages back. A single forward
pass has no architectural verifier; coherence is a side effect of
training, not an enforced invariant.

Newer systems narrow the gap. Reasoning-trained models, agent
loops that ask a separate model to critique a draft before it
ships, and inference-time techniques like sampling several
completions and comparing them all reduce inconsistency at the
cost of compute and orchestration complexity. None is a guarantee,
and most legal-AI tools currently expose only a fraction of these
capabilities to the user. Whether your tool runs a self-check pass
on its own output --- and what kinds of inconsistency it is
trained or tuned to catch --- is a useful question to put to your
vendor.

The remedy at the lawyer's end remains unsexy: read the draft
slowly, with the eye of a reviewer who is reading for internal
contradiction rather than for factual correctness. Domain
expertise unmasks the subtle ones. A junior associate who
understands what indemnification typically looks like will catch a
clause that flips the allocation; a reviewer who is only checking
citations will not.

\section*{Context: is the output appropriate for the situation?}

A structurally valid, internally coherent draft can still answer
the wrong question. A client memo written in motion-brief register
is wrong even if every sentence is true. A research summary
covering all federal law is wrong if the matter is in a state
court of limited jurisdiction. A demand letter written in
aggressive litigation prose is wrong if the client wants to
preserve the commercial relationship. A brief written for a
generalist trial judge is wrong if it is going to a specialty
court that does not need the doctrinal background.

The mechanism is the gap between what the lawyer knows about the
situation and what the prompt encodes. The model has access only
to the tokens in its context window. Anything the lawyer knows
about the audience, the jurisdiction, the tone, the strategic
posture, or the desired length that is not written into the prompt
does not exist for the model. System prompts and house templates
shape defaults but defaults are defaults, not the situation.

Context failures are easy to recognize and harder to remedy. The
reviewer can usually point at the answer and name what is wrong
with it (``too long,'' ``wrong audience,'' ``misses the actual
question'', ``wrong tone for this client''). The fix is to give
the model more of the situation, in writing, before the next
draft. The asymmetry between the ease of recognition and the
difficulty of fix is the diagnostic signature of a context failure
in legal work as elsewhere.

\section*{Groundedness: do the legal claims connect to actual law?}

This is the rung the profession has learned to fear, and rightly.
Hallucinated cases. Hallucinated holdings. Misquoted statutes.
Made-up regulatory thresholds. Mis-summarized procedural rules.
The Magesh study [2] showed that even retrieval-augmented legal
tools, designed to ground their outputs in real authority, fail at
meaningful rates. Retrieval reduces hallucination; it does not
eliminate it. Understanding why helps set the right expectations
for any retrieval-augmented tool.

When a tool retrieves the text of a case and hands it to the
language model, the model does not transcribe the holding. It
produces the most plausible continuation of the prompt given the
retrieved text, which is usually close to faithful but is
sometimes a fluent paraphrase that subtly alters emphasis,
scope, or qualification. The model has a prior, baked in during
training, about what a case is likely to say based on its name,
jurisdiction, and doctrinal neighbourhood. When that prior
conflicts with the retrieved text, the model does not always
defer to the text; it sometimes resolves the conflict by
producing the answer its prior predicts. The result is an output
that cites a real case, points to a real source, and still
misrepresents the holding --- not fabrication in the Mata sense,
but misreading in the way a rushed human reader might misread,
at scale and without the fatigue that would signal the risk.
The model still has to read the retrieved authority
faithfully, defer to it when its parametric prior conflicts,
and reconcile contradictions among sources. Each of those steps
is a place where groundedness can fail even when the citation is
real.

We use \emph{groundedness} for the property the output either has
or does not. \emph{Grounding}, in the broader literature, names
the underlying problem: how a system that has only seen text can
come to refer reliably to facts and authorities outside the text.
The hallucination literature splits the failures further: the
output may contradict the world (the cited case does not exist) or
contradict a source the model was given (the cited case exists but
does not say what the model summarized it as saying). Both are
groundedness failures in the sense used here, and both are
catchable only from outside the AI's output.

The difficulty is compounded by a calibration problem specific to
legal work. The model's expressed confidence --- the apparent
authority of its prose --- gives no reliable signal of accuracy.
A hallucinated holding arrives in the same authoritative register
as a verified one. A misquoted statutory threshold sounds as
precise as the actual text. Legal writing is authoritative in
tone by convention; the model has absorbed that convention and
applies it uniformly, whether the underlying claim is reliable or
fabricated. There is no surface signal that distinguishes a
misread case from a correctly read one. The implication for
practice is that the plausibility of a summary cannot substitute
for its verification. A persuasive-sounding statement of a
holding is not evidence that the statement is accurate.

The discipline for legal practice is the one the bar has long
preached: verify any cited authority that matters. Use legal
research tools that produce checkable citations and check them.
Treat any factual or legal claim by the AI as a hypothesis until a
human or a citator has confirmed it. The careful firm builds
review workflows in which groundedness is catchable by
construction --- citation auditing, source linking, mandatory
KeyCite or Shepard's verification on every cite the AI produces.

\section*{Environment: did the prompt encode what you know?}

The fifth rung sits at the human end of the ladder and is the
hardest to see because, by construction, it is invisible. An
environment failure occurs when the answer is wrong because the
prompt did not contain enough of the situation for the question to
be answerable. The lawyer knows things --- about the client's
risk tolerance, the partner's drafting preferences, the local
court's standing orders, the opposing counsel's history, the
matter's strategic posture, the firm's house style, the events of
yesterday afternoon's call --- that the model does not. If the
lawyer does not write those things into the prompt, the model
cannot use them. The output can be fluent, internally consistent,
properly grounded, and still wrong because the question itself was
underspecified.

The signature is distinctive. The reviewer reads the draft and
feels frustrated. The draft is, on its own terms, fine. What is
wrong is that it does not address the real situation. The fix is
rarely to ask the model again. The fix is to write down,
explicitly, the parts of the environment the model needed to see.

This is the rung most legal AI policies underweight. Citation
checkers catch some structure and groundedness problems. Tone and
audience checks catch some context problems. But environment
errors --- the demand letter that misreads the client's strategic
posture, the research memo that surveys settled law because the
prompt did not say which threshold issue had already been decided
--- escape every automated tool. The lawyer's environment is not
on the firm's intranet, in the matter management system, or
attached to the email that arrived this morning. It lives in the
lawyer's head. Importing it into the prompt is the lawyer's job.

\section*{Beyond the five rungs}

The five rungs map single-output failures --- everything
diagnosable when a specific draft is in front of you. Six
additional failure classes exist alongside them; they escape
the rung-by-rung frame and require different detection and
remediation. The ladder addresses only the five rungs. Knowing
what lies beyond them before walking makes the ladder's scope
explicit and its limits visible.

The first is the \emph{normative} failure. An output can be
structurally correct, internally coherent, situationally
appropriate, grounded in real authority, and still wrong ---
because it reflects a value judgement the lawyer or client did
not authorise. An AI that applies a statute mechanically to a
set of facts may produce a conclusion that is legally defensible
but strategically, ethically, or commercially inadvisable for
this client in this matter. The five rungs ask whether the answer
is correct. The normative question asks whether the answer is
right, and right is not the same thing as correct. That
judgement belongs to the lawyer, not the tool.

The second is the \emph{adversarial} failure. In document review
and contract analysis, the AI processes materials that an
opposing party or third party supplied. A hostile party who
understands how the tool works can craft a document --- a
contract redline, a discovery production, a third-party
submission --- whose text is designed to redirect the model's
analysis or override its instructions. This is roughly what
security researchers call prompt injection: malicious
instructions embedded in content the model is asked to process.
The ladder will identify the output as a context or groundedness
failure, but the diagnosis alone does not reveal that the failure
was induced deliberately. Firms using AI for document review
in adversarial contexts should treat this as a live risk, not a
theoretical one.

The third is the \emph{drift} failure. The model's parametric
knowledge of case law, statutes, and regulatory guidance was
fixed at its training cutoff. Legal practice surfaces this
regularly: a case that was good law when the tool last trained
may have been distinguished, reversed, or superseded by the time
the memo reaches the client. Retrieval-augmented tools partially
address this by fetching current authority, but only for the
specific queries they run; the model's background assumptions
about the doctrinal landscape remain frozen. For rapidly evolving
areas --- regulatory rulemaking, recent circuit splits, new
statutory amendments --- the gap between what the model ``knows''
and what the law currently says is not a grounding failure in the
fabrication sense. It is a staleness failure that passes every
rung of the ladder until someone with current doctrinal knowledge
reads the output. Currency checking, in areas where it matters,
remains a human responsibility.

The fourth is the \emph{complacency} failure. The duty of
competence is not only a duty to establish a review process; it
is a duty to sustain the quality of that review over time. A
firm whose attorney oversight becomes perfunctory after the tool
has produced a long run of correct outputs has not satisfied the
duty it described to the ethics committee. Automation complacency
is well-documented in human-factors research: a reviewer who has
processed many correct outputs processes the next one less
vigilantly. The bar records whether review happened, not how
carefully. The tool's failure rate may be constant; the effective
error rate seen by clients rises as vigilance decays.

The fifth is the \emph{compositional} failure. Legal analysis
characteristically requires chaining multiple inferences:
determining standing, then jurisdiction, then the applicable
standard of review, then whether the facts satisfy each element
of the test. When a model works through a multi-step analysis,
errors compound across the chain. A model that is highly reliable
on each individual step may produce a wrong final conclusion
because a subtle error in an early step --- a mischaracterization
of the standard, a slightly wrong reading of the procedural
posture --- propagates through the downstream reasoning. No single
step fails the five-rung ladder in a way that a careful reading
of that step alone would catch; the failure is a property of the
chain, not of any individual link. Multi-step doctrinal analysis
deserves verification not only at the conclusion but at each
material intermediate step.

The sixth is the \emph{coverage} failure. A legal AI tool is
shaped not only by what it has seen but by what its fine-tuning
and evaluation emphasized. Tools trained and benchmarked on broad
legal research tasks may have attenuated, lower-quality performance
in specific niches --- specialty courts, narrow regulatory domains,
emerging statutory areas, non-majority jurisdictions --- not
because the relevant authority is absent from training data but
because the model's behavior was shaped toward the modal legal
task, not toward the specific practice area. A tool that performs
impressively on mainstream commercial litigation research may be
unreliable in the particular procedural context, doctrinal niche,
or jurisdiction of a given matter. A coverage failure can present
on the ladder as a grounding failure: the model states holdings,
statutory elements, or procedural requirements with confidence,
and they turn out to be wrong for the specific practice area. The
ladder correctly names the symptom as a grounding failure; it does
not identify the cause as a training-coverage mismatch. Applying
the grounding remedy --- retrieval, citation checking, source
verification --- addresses individual errors but not the
structural problem. The structural remedy is targeted evaluation
on representative examples from the specific practice area before
deploying the tool in that context, rather than relying on general
benchmark performance.

These six failure classes require different responses than the
ladder provides: professional-conduct governance for normative
failures; security protocols for adversarial failures;
currency-checking workflows for drift; sustained oversight
culture for complacency; intermediate-step verification for
compositional failures; practice-area evaluation before
deployment for coverage failures. The ladder below does not
address them. It is a procedure for the five-rung failures ---
the class that shows up most often in practice and that yields
most directly to prompt discipline, retrieval, and reviewer
skill.

\section*{Walking the ladder}

The five rungs form a diagnostic ladder ordered from what the AI
can most reliably handle to what only the lawyer can supply.
With the full failure landscape in view, the ladder's scope is
clear: it is a procedure for the five-rung failures. A senior
associate sends you a draft summary judgment brief generated with
AI assistance. You walk the rungs.

\emph{Structure} passes: the brief is properly formatted, the
citations are Bluebooked, the table of authorities is generated.
\emph{Meaning} passes: the argument is internally consistent.
\emph{Context} fails: the brief is written in the aggressive
register the partner reserves for trial-bound matters, but this is
a settlement-track case in which the client wants to look
reasonable to the bench. You revise the prompt to specify the
strategic posture and the desired register. The brief comes back
softer.

\emph{Groundedness} fails on cite-check: one of the supporting
cases has been distinguished in a more recent appellate decision
and another, on closer reading, does not actually stand for the
proposition the brief cites it for. You ask for the AI's source
links, audit them, and instruct the next draft to use only cases
the citator confirms as currently good law for the proposition
asserted.

\emph{Environment} fails last. The partner's strategy is to lead
with the dispositive jurisdictional argument and reserve the
substantive arguments for a later motion if the jurisdictional
play fails. The AI did not know that. The brief covers all four
issues with equal weight. You add the strategic priority to the
prompt --- ``lead with jurisdictional, hold the substantive
arguments in reserve'' --- and the next draft reflects it.

The walk located each failure at the rung where the remedy lives.
Not one of these failures was a hallucination in the headline
sense. Each would have escaped a workflow that checked structure
and citations and stopped.

\section*{What this means for firms}

The ladder suggests four concrete things for firms managing AI in
legal practice.

\emph{First}, review workflows should walk all five rungs, not
just structure and groundedness. The middle rungs --- meaning and
context --- are the ones most often missed by automated tooling
and most often cited by partners as ``something feels off about
this draft.'' A structured rubric helps.

\emph{Second}, the rungs map to different parts of the playbook.
Structural problems call for better tooling (citation linters,
contract validators). Groundedness problems call for better
retrieval (verified-citation tools, mandatory citator checks).
Context and environment problems call for better prompt-and-brief
discipline by the lawyer in front of the tool. A firm that wants
to reduce groundedness errors should invest in retrieval and
citation auditing. A firm that wants to reduce environment errors
should invest in training, templates, and reviewer practice.
Knowing which rung you are on tells you which budget line to
charge.

\emph{Third}, the lawyer's continuing duty of competence and
supervision under the relevant rules of professional conduct ---
addressed for AI tools in ABA Formal Opinion 512 [3] --- is in
practice a duty to walk the higher rungs. The AI handles the
lower rungs well. Verification of grounding is now table stakes.
Context and environment are the rungs at which the lawyer's
judgment is most clearly indispensable, and they are the rungs
that no current tool can climb on the lawyer's behalf.

\emph{Fourth}, firms should account for verification cost
explicitly. AI reduces the time to produce a first draft; it
does not reduce the time to verify that draft at the same rate.
A business case built on drafting-time savings alone
systematically undercounts the verification burden the five-rung
framework makes visible. The bottleneck shifts from drafting to
careful review, and workflows that treat review as an afterthought
will find that generation savings are consumed by downstream costs.
The ladder is useful here too: it tells review teams where to
concentrate, because the failure modes on each rung require
different skills and different checks.

The ladder is not a theory of legal AI. It is a procedure for the
moment in which a draft arrives, looks plausible, and is not quite
right, and someone needs to say --- in plain English, before the
filing deadline --- what kind of wrong it is.

\section*{References}

\noindent
[1] \emph{Mata v.~Avianca, Inc.}, No.~22-cv-1461, 2023 WL 4114965
(S.D.N.Y. June 22, 2023) (sanctioning attorneys for filing a brief
with fabricated case citations generated by ChatGPT).

\noindent
[2] V.~Magesh, F.~Surani, M.~Dahl, M.~Suzgun, C.~D.~Manning, and
D.~E.~Ho, ``Hallucination-Free? Assessing the Reliability of
Leading AI Legal Research Tools,'' Stanford RegLab / HAI working
paper, 2024.

\noindent
[3] American Bar Association Standing Committee on Ethics and
Professional Responsibility, \emph{Formal Opinion 512: Generative
Artificial Intelligence Tools}, July 29, 2024.

\bigskip

\noindent
\textit{Lodewijk Bonebakker is the author of \emph{Large Language
Models as Engineered Systems}, a forthcoming primer for engineers
on the systems, statistical, and human-factors dimensions of
modern LLMs. He can be reached at
\texttt{focus2flow@bonebakker.org}. The author is not a practicing
attorney; this article is intended as a framework for thinking
about AI errors in legal work, not as legal advice.}

\end{document}
